\documentclass[12pt]{article}
\usepackage{amsmath, amssymb, amsthm}
\usepackage{geometry}
\usepackage{algorithm2e}
\usepackage{booktabs}
\geometry{margin=1in}

\title{Homework: Recurrence Equations}
\author{Algorithm Design and Analysis}
\date{}

\begin{document}
\maketitle

\section{Problem Statement}

Solve the following recurrence equations and express the solution in $\Theta$-notation.

\begin{enumerate}
  \item $T(n) = 2T(n/2) + n$
  \item $T(n) = T(n/2) + 1$
  \item $T(n) = 4T(n/2) + n^2$
  \item $T(n) = 3T(n/3) + n \log n$
  \item $T(n) = T(n-1) + n$
  \item $T(n) = 2T(n-1) + 1$
\end{enumerate}

\section{Background: The Master Theorem}

Given a recurrence of the form
\[
  T(n) = aT(n/b) + f(n), \quad a \geq 1,\ b > 1,
\]
let $c^* = \log_b a$. Then:

\begin{description}
  \item[Case 1] If $f(n) = O(n^{c^*-\varepsilon})$ for some $\varepsilon > 0$, then $T(n) = \Theta(n^{c^*})$.
  \item[Case 2] If $f(n) = \Theta(n^{c^*} \log^k n)$ for some $k \geq 0$, then $T(n) = \Theta(n^{c^*} \log^{k+1} n)$.
  \item[Case 3] If $f(n) = \Omega(n^{c^*+\varepsilon})$ and the regularity condition holds, then $T(n) = \Theta(f(n))$.
\end{description}

\section{Solutions}

\subsection*{(1) $T(n) = 2T(n/2) + n$}
$a=2, b=2, c^*=\log_2 2=1$. Here $f(n)=n=\Theta(n^1)$, so Case~2 applies ($k=0$).
\[
  \boxed{T(n) = \Theta(n \log n)}
\]

\subsection*{(2) $T(n) = T(n/2) + 1$}
$a=1, b=2, c^*=0$. Here $f(n)=1=\Theta(n^0)$, so Case~2 applies.
\[
  \boxed{T(n) = \Theta(\log n)}
\]

\subsection*{(3) $T(n) = 4T(n/2) + n^2$}
$a=4, b=2, c^*=2$. Here $f(n)=n^2=\Theta(n^2)$, so Case~2 applies.
\[
  \boxed{T(n) = \Theta(n^2 \log n)}
\]

\subsection*{(4) $T(n) = 3T(n/3) + n\log n$}
$a=3, b=3, c^*=1$. Here $f(n)=n\log n=\Theta(n^1\log n)$, so Case~2 applies ($k=1$).
\[
  \boxed{T(n) = \Theta(n \log^2 n)}
\]

\subsection*{(5) $T(n) = T(n-1) + n$ (substitution method)}
Unrolling: $T(n) = n + (n-1) + \cdots + 1 = \frac{n(n+1)}{2}$.
\[
  \boxed{T(n) = \Theta(n^2)}
\]

\subsection*{(6) $T(n) = 2T(n-1) + 1$ (substitution method)}
Unrolling: $T(n) = 2^n T(0) + 2^{n-1} + \cdots + 1 = \Theta(2^n)$.
\[
  \boxed{T(n) = \Theta(2^n)}
\]

\section{Recursion Tree Analysis (Problem 1)}

For $T(n) = 2T(n/2) + n$, the recursion tree has $\log_2 n$ levels.
At level $i$ there are $2^i$ nodes each contributing $n/2^i$ work.
Total work per level: $n$. Total over all levels:
\[
  T(n) = n \cdot \log_2 n = \Theta(n \log n).
\]

\end{document}
